\chapter{JUSTIFICATIVA}
\label{capitulo-justificativa-proposta}

A mobilidade independente em recintos fechados e logradouros públicos constitui um dos fatores mais determinantes para a autonomia e a inclusão social de pessoas com baixa visão. No cotidiano, a presença de obstáculos aéreos e barreiras fora do alcance do solo representa um perigo constante de impacto corporal e traumatismo craniofacial. Nesse contexto, recursos de tecnologia assistiva têm papel fundamental para ampliar a segurança pessoal e restabelecer a confiança no deslocamento autônomo.

\section{Motivação do Projeto e Resultados Esperados}
\label{sec:motivacao-esperada-proposta}

Este trabalho existe para suprir a zona cega deixada pelos instrumentos convencionais de mobilidade, como a bengala longa, propondo um dispositivo complementar fixado ao punho que monitora o espaço superior e frontal. Espera-se com o protótipo fornecer alertas sensoriais discretos via resposta vibrotátil (feedback háptico), permitindo que o usuário identifique a proximidade de obstáculos sem comprometer a sua percepção auditiva do ambiente --- canal sensorial imprescindível para a ecolocalização e atenção a riscos no trânsito.

A seleção do microcontrolador ESP32 justifica-se por sua capacidade de processamento dual-core de 32 bits, dimensões compactas, suporte nativo a modulação por largura de pulso (PWM) por hardware e excelente relação custo-benefício. Ao empregar componentes comerciais prontos para uso (\textit{COTS}), o projeto visa demonstrar a viabilidade técnica de uma tecnologia assistiva de código aberto e baixo custo, cujo valor de montagem seja acessível à realidade socioeconômica brasileira.

\section{Relevância para a Realidade Local e Comunitária}
\label{sec:relevancia-local-proposta}

A relevância desta proposta para a comunidade reside na democratização de soluções assistivas. No mercado atual, dispositivos eletrônicos comerciais de auxílio à navegação são quase exclusivamente importados, protegidos por patentes e comercializados a valores proibitivos para a maior parte da população. Desenvolver um protótipo aberto e funcional nos laboratórios do IFPR Campus Pinhais aproxima o ensino público e a pesquisa aplicada das necessidades reais de acessibilidade do município e região.

\section{Implicações da Não Realização da Ação}
\label{sec:nao-realizacao-proposta}

Caso a ação de extensão não seja realizada, perde-se a oportunidade de validar uma alternativa tecnológica inclusiva, aberta e de baixo custo desenvolvida no ambiente acadêmico do IFPR. Além disso, deixa-se de promover uma valiosa aproximação entre o conhecimento científico do curso de Ciência da Computação e a sociedade, mantendo os estudantes afastados da experiência prática de solucionar problemas reais de acessibilidade física e cidadania.
